\section{Câu hỏi trước phương pháp}
\label{sec:ch06-cau-hoi-truoc-phuong-phap}

\index{attribution}%
Mục cuối của chương~\ref{ch:attribution-theo-gradient} đã tách bảo đảm hình
thức khỏi kết luận về độ trung thực. Sự tách ấy cũng áp dụng khi so các họ
phương pháp. Taimeskhanov và Garreau bắt đầu từ một định nghĩa tối giản lấy
từ tài liệu trước.

\begin{definition}[Attribution đặc trưng]
\label{def:ch06-attribution}
Một attribution tại đầu vào $x$ cho mô hình $f$ là một vector
$\varphi(x,f)$, trong đó thành phần $\varphi_j(x,f)$ được diễn giải là đóng
góp của đặc trưng $j$~\autocite{p17firstprinciples}.
\end{definition}

Định nghĩa~\ref{def:ch06-attribution} chưa nói điểm số được lấy bằng phép so
sánh, phép trung bình hay đạo hàm nào.

LIME, SHAP và Integrated Gradients đã cho ba câu trả lời khác nhau cho khoảng
trống đó. LIME chọn lân cận và trọng số cho các mẫu đã nhiễu. SHAP chọn giá trị
của liên minh và một phân phối nền. Integrated Gradients chọn baseline cùng
đường đi tới đầu vào. Vì vậy, việc hai phương pháp trả các điểm số khác nhau
không tự chứng minh phương pháp nào sai. Chúng có thể chỉ đang trả lời
những câu hỏi \tn{phản thực}{counterfactual} khác nhau.

Bài báo của Taimeskhanov và Garreau đặt câu hỏi ở mức trừu tượng hơn ba
phương pháp ấy. Thay vì đặt
điều kiện cho một công thức đã có, bài báo hỏi hai câu theo thứ tự. Câu thứ
nhất: với lớp mô hình đơn giản nhất có thể, attribution nào là đáng tin? Câu
thứ hai: từ attribution của những mô hình đơn giản ấy, ghép thành attribution
của một mô hình $f$ bất kỳ bằng cách nào~\autocite{p17firstprinciples}? Cách
đặt câu hỏi này sẽ buộc lựa chọn ẩn trong mỗi phép tính lộ ra.
